\addtocontents{toc}{\protect\Needspace{3\baselineskip}}
\originalchapter{闭超弦类型与平行 D 膜上的谱}\label{ch:basic-types-branes}
对应 8.251 第 23--24 讲. 本章把上一章的左右扇区组合成十维闭超弦, 再将开弦边界分成 NN、DD、ND 类型. 五种超弦及其对偶的历史概况保留为有条件的研究结果, 不当作已经从自由振子推导出的全部统一理论.

\originalsection{Type II 的四个扇区}
闭超弦有两组独立世界面费米模式, 两边各自作 GSO 投影, 并要求物理层匹配. 扇区为 NS--NS、NS--R、R--NS、R--R. 前者以及 R--R 给时空玻色态, 两个混合扇区给费米态. 左右 R 零模选同手征得 Type IIB, 选相反手征得 Type IIA.

无质量 NS--NS 为 $8_v\otimes8_v$, 分成对称无迹 35、反对称 28、迹 1, 对应十维引力子、二形式与膨胀子. 两个混合扇区各有 $8\times8=64$ 态, 共 128 个费米极化. R--R 的 64 态依手征给不同的形式场:
\[
 \text{IIA: }C_1,C_3,\qquad
 \text{IIB: }C_0,C_2,C_4\ (F_5=*F_5).
\]
IIA 数目为 $\binom81+\binom83=8+56=64$; IIB 为 $1+28+\frac12\binom84=64$. 自对偶约束使四形式的极化数减半, 不能多算一份. NS--NS 与 R--R 共 128 个玻色极化, 与两个混合扇区数目相等.

\originalsection{形式势与膜的电磁维数}
一个 $p$ 空间维膜可通过 $\int_{W_{p+1}}C_{p+1}$ 耦合到 $p+1$ 形式势, 成为该势的电源. 在 $D$ 维, 若场强为 $p+2$ 形式, 磁对偶场强为 $D-p-2$ 形式, 相应磁膜空间维数为 $D-p-4$. IIA 中 D0、D2 分别为 $C_1,C_3$ 的电源; IIB 中 D1、D3 分别耦合 $C_2,C_4$, D7 是 $C_0$ 的磁对偶源. 十维 D3 的磁对偶仍是 D3, 与 $F_5$ 自对偶相容.

D3 在其四维世界体积内可以“充满空间”, 在完整十维目标空间却有六个法向方向. 原第 23 讲的“space filling”应在相应低维语境理解, 不能说十维 D3 是充满九个空间方向的 D9.

\originalsection{异质弦的左右中央荷}
异质弦左边采用玻色构造, 右边采用超弦构造. 十维非紧致坐标在左边提供中央荷 10, 还需要中央荷 16 的内部体系才能达到玻色条件 26; 右边则用十维超坐标和超共形结构. 可以用 16 个紧致内部玻色自由度描述左边额外体系. 因此“$26-10=16$”是中央荷配置, 不意味着可见物理时空左右各有不同维数.

在标准偶自对偶 16 维内部格构造中, 两个选择给 $E_8\times E_8$ 与 $\operatorname{Spin}(32)/\mathbb Z_2$ 型异质弦. 此处引用格与模不变性的一致性结果, 不把一个维数相加等式当作对这两个群的分类证明. 若内部格向量为无量纲 $P$, 其左层贡献为 $P^2/2$, 质量条件写成
\[
 \frac{\alpha'M^2}{4}=N_L-1+\frac{P^2}{2}
 =N_R-a_R,
 \quad a_R=\tfrac12\ \text{(NS)},\quad0\ \text{(R)}.
\]
根向量 $P^2=2$ 的左边基态可产生无质量规范极化. 这是看到内部自由度如何成为低维规范场的入口, 完整紧致化谱需另指定格、Wilson 线和背景.

\originalsection{五种超弦的概况与边界}
通常称为五种十维微扰超弦的是 IIA、IIB、两种异质弦和带开弦及非定向投影的 Type I. 标准 Type I 的规范群与异常抵消要求相联系, 不能只从“有 32 个标签”推出一致性. 各种强弱耦合对偶和 M 理论极限把这些构造联系起来, 但不存在本卷已推导的统一非微扰作用量. 原课的历史示意图在此转为这个明确的范围说明, 避免将推测中的统一表述写成一套已穷尽的定理.

\originalsection{单膜上的 Neumann 与 Dirichlet 模式}
设 D$p$ 膜沿 $x^0,\ldots,x^p$, 光锥取膜内 $x^0,x^1$, 因而此规范要求 $p\geq1$. 独立横向有 $p-1$ 个 NN 方向 $i=2,\ldots,p$, 及 $D-p-1$ 个 DD 方向 $a=p+1,\ldots,D-1$. NN 包含连续零模动量, DD 不含独立法向动量零模, 但仍含正范数振动模式.

同膜 DD 坐标可取实展开
\[
 X^a=y^a+\sqrt{2\alpha'}\sum_{n\ne0}\frac{\alpha_n^a}{n}
 \ee^{-\ii n\tau}\sin n\sigma,
 \qquad\alpha_{-n}^a=(\alpha_n^a)^\dagger.
\]
这里不能在前面再加 NN 展开的 $\ii$: $n\mapsto-n$ 时 $\sin n\sigma$ 与 $1/n$ 两个负号相消, 上式才厄米. 直接计算 $X^a(0)=X^a(\pi)=y^a$ 且求正则核可核对振子仍满足 $[\alpha_m^a,\alpha_n^b]=m\delta^{ab}\delta_{m+n,0}$.

临界玻色开弦仍有 $\alpha'M^2=N-1$, 这里 $M$ 是膜上场的质量, 连续动量标签只在膜内. $N=1$ 的 NN 态给膜上向量的 $p-1$ 个物理极化, DD 态给 $D-p-1$ 个标量. 标量表示膜的法向位移, 不是能在法向自由传播的额外零模粒子. D0 需另用世界线规范, 没有传播的世界线光子; 其法向位置变量仍存在.

\originalsection{平行膜间的伸长质量}
两膜相距向量 $b^a=y_2^a-y_1^a$. DD 展开改为
\[
 X^a=y_1^a+\frac{b^a}{\pi}\sigma+\text{正弦振动},
\]
线性项给经典伸长能 $T_s|b|$, 由约束得
\[
 M^2=(T_s|b|)^2+\frac{N-1}{\alpha'}
 =\frac{|b|^2}{4\pi^2\alpha'^2}+\frac{N-1}{\alpha'}.
\]
对 GSO 后超弦, 将截距与允许 $N$ 换成对应扇区, 无质量层的延伸态质量是 $|b|/(2\pi\alpha')$. 两膜重合时这些态变轻, 与非阿贝尔规范增强的 Higgs 解释相容.

有定向弦的端点标签为 $[i,j]$, $[i,j]$ 与 $[j,i]$ 一般不同. 弦拼接 $[i,j][j,k]\to[i,k]$ 对应矩阵乘法; $N$ 张重合膜的标签为 $N\times N$ 矩阵, 导向 $U(N)$ 规范变量. 单靠计数尚不证明相互作用作用量, 第二编从 DBI 与低能展开作进一步说明.

\originalsection{混合边界与半整数玻色模}
若一个方向在一端为 Neumann、另一端 Dirichlet, 可取 $\cos r\sigma$ 或 $\sin r\sigma$, $r=n+1/2$. 例如 $X'(0)=0,X(\pi)=0$ 给 $\cos r\pi=0$, 因而半整数频率, 无位置或动量零模. 一个 ND 玻色方向的有限零点能是 $\frac12\zeta(-1,1/2)=+1/48$, 比 NN/DD 的 $-1/24$ 增加 $1/16$. 对有 $\nu$ 个 ND 方向的临界玻色弦,
\[
 \alpha'M^2=N-1+\frac{\nu}{16}
 +\frac{|b_{\rm DD}|^2}{4\pi^2\alpha'}.
\]
ND 与 DD 伸长来源不同, 不能把所有垂直坐标都当伸长零模. 第 26 讲复习题将具体检查一个 ND 与 23 个整数模的例子.
